problem:  
Let \( D \subset \mathbb{C}^n \) be an irreducible bounded symmetric domain realized via its Harish–Chandra embedding with Bergman metric \( g_B \), and let \( S(D) \) denote its Shilov boundary. For a smoothly bounded, strictly pseudoconvex domain \( \Omega \subset \mathbb{C}^n \) with Bergman metric \( g_\Omega \), assume that near some boundary point \( p \in \partial\Omega \), the CR invariants of \( \partial\Omega \) match—up to order \( k \ge 3 \)—those of a smooth CR submanifold \( N \subset S(D) \); concretely, the Levi form, the Webster scalar curvature, and all covariant derivatives of order up to \( k \) of the Tanaka–Webster curvature tensor agree under a local CR diffeomorphism 
\[
\Phi: (\partial\Omega, p) \to (N, q) .
\]
  
**Question:**  
Does there exist a critical finite order \( k_* \), depending only on the symmetric domain type of \( D \), such that CR curvature matching up to order \( k_* \) forces \( \Omega \) to be locally biholomorphic to a domain in \( D \)?  
Equivalently, is finite–order boundary equivalence to a symmetric Shilov model sufficient to determine the interior Bergman geometry up to holomorphic isometry?  

Why it is a "good" problem:  
This problem asks for a *finite-order* boundary-to-interior rigidity principle for Bergman metrics, bridging Fefferman–Graham CR invariant theory and the intrinsic Kähler symmetry of bounded symmetric domains.  

While classical CR results (Chern–Moser, Webster, Baouendi–Ebenfelt–Rothschild) establish that **full jet** equivalence at all orders enforces local biholomorphism, the finite‑order threshold for such rigidity remains unknown even for the unit ball. Determining such a universal \( k_* \) would quantify how much boundary data one must know to recover an interior symmetric structure—analogous to finite-jet determinacy problems in differential geometry and analysis.  

The problem lies at the intersection of:  

- **CR geometry** (finite‑jet equivalence, Tanaka–Webster connections),  
- **Several complex variables** (Bergman kernel asymptotics, boundary regularity), and  
- **Representation theory** (invariance under isotropy actions of \( \mathrm{Aut}(D) \)).  

It is conceptually simple—“Is there a finite order of boundary curvature matching that forces symmetry?”—but solving it would require breakthroughs joining CR invariant theory with metric analysis on Kähler–Einstein domains. This finite‑order rigidity threshold problem is open, likely delicate, and would deepen our understanding of the analytic continuation of symmetric geometric data from boundary to interior.  

Hence, it exemplifies the “narrow entrance, profound depth” aesthetic and would constitute a significant advance in the analytic–geometric theory of Hermitian symmetric spaces.